\section{Product exponents and local-curve conventions}\label{sec:integrality}

The factorization of Lemma~\ref{lem:formal} is rational. The
compatible series $f_0(Q)^{1/2}$ shows that the formal axioms alone
cannot imply integral exponents. For projective Calabi--Yau
threefolds, integer exponents in this particular product expansion
are the content of \cite[Conjecture 1.4]{Pardon}.

The local series and the elementary series of Ionel--Parker
\cite{IP} arise from different torus actions. The distinction is
visible in an exact degree-two calculation.
Let $F_g(q,u)=1+\sum_{d\ge1}\GW(x_{g,d})q^d$, where $q$ records
curve degree and $u$ is the GW variable. Pardon's
Proposition 4.8 identifies this with the residue theory for scalar
fiberwise scaling. Let $E_g(q,u)$ be the exponential of the
connected elementary series of \cite[Proposition 3.3]{IP}, given by
the antidiagonal specialization
\[
 E_g(q,u)=\sum_\mu q^{|\mu|}
    \prod_{\square\in\mu}(2\sin(h(\square)u/2))^{2g-2}.
\]
Both use balanced level $(g-1,g-1)$. Put $s=2\sin(u/2)$.
The scalar formula of Bryan--Pandharipande
\cite[Section 8]{BP}, and the two partitions of two in
the hook formula, respectively give
\[
 [q]F_g=[q]E_g=s^{2g-2},
\]
\[
 [q^2]F_g=s^{4g-4}\bigl((4-2s)^{g-1}+(4+2s)^{g-1}\bigr),
 \qquad [q^2]E_g=2s^{4g-4}(4-s^2)^{g-1}.
\]
In particular
\begin{equation}\label{eq:discrepancy}
 [q^2](\log F_2-\log E_2)=2s^6,\qquad
 F_2(q,u)=E_2(q,u)E_4(q^2,u)^2\pmod{q^3}.
\end{equation}
Thus direct identification fails, although this first correction
has an integer exponent. More generally, for $g\ge2$ binomial
expansion gives
\[
 F_g(q,u)=E_g(q,u)\prod_{j=1}^{g-1}
 E_{2g-1+j}(q^2,u)^{a_{g,j}}\pmod{q^3},
\]
\[
 a_{g,j}=2\,4^{g-1-j}
 \left(\binom{g-1}{2j}-(-1)^j\binom{g-1}{j}\right)\in\ZZ,
\]
with out-of-range binomial coefficients zero. The degree-two
discrepancy vanishes for $g=0,1$.

At degree three the comparison has a fractional correction. For a
primitive class $\alpha$, write $c_{\alpha,g}$ for the integer
exponents in the elementary expansion of \cite[Theorem 8.1]{IP},
and $e_{\alpha,g}$ for the scalar-local exponents of
\eqref{eq:factor}. The companion paper \cite{WuebbenIntegrality}
proves
\[
 e_{3\alpha,10}=c_{3\alpha,10}-\frac{4720}{3}c_{\alpha,3}.
\]
On the resolution $Y_+$ of Section~\ref{sec:geometry}, the primitive
class $\alpha=3b_2+4b_3$ has $c_{\alpha,3}=-14$. Consequently
$e_{3\alpha,10}\in\ZZ+2/3$. This disproves the integer-exponent
assertion of \cite[Conjecture 1.4]{Pardon} for its scalar generators,
even with torsion-free second homology. It does not contradict the
integer elementary expansion or the distinct integral geometric
generators in \cite[Theorem 1.2]{Pardon}.

Corollary~\ref{cor:integral} describes how the rational exponents
specialize. Its integrality consequence applies whenever the source
exponents are integral; the counterexample shows that this condition
cannot be inferred from projectivity and the Calabi--Yau hypothesis.
None of the support arguments or universal coefficient calculations
above requires it.
